\chapter{Investigaciones previas relacionadas}
\label{chap:investigaciones-previas}

\insertminitoc
\parindent1.5em

\section{Estado actual de la investigación y antecedentes}

La investigación sobre la infancia en la era digital ha dado un giro crítico tras los eventos globales de la última década, pasando de estudiar el simple acceso a la tecnología a analizar la profundidad de la inmersión virtual. Según el análisis de \cite{limone_psychological_2021}, la pandemia de COVID-19 no solo incrementó exponencialmente el tiempo de pantalla, sino que forzó una reconfiguración neuropsicológica en los menores. El aislamiento social convirtió a los dispositivos móviles en el único "cordón umbilical" con la realidad externa y con el grupo de pares, derivando en una dependencia digital que la comunidad científica asocia hoy con alteraciones severas en la regulación emocional. Esta visión clínica es complementada por \cite{holmarsdottir_integrative_2025}, quienes proponen un enfoque ecológico para entender el bienestar digital subjetivo. Su marco teórico sugiere que las investigaciones deben trascender el simple conteo de horas frente a la pantalla y centrarse en cómo el niño percibe sus interacciones en los dominios superpuestos de la familia, la educación y el ocio digital.\\



Para los niños en la primera infancia (0 a 5 años), el impacto es biológicamente más sensible debido a la acelerada plasticidad cerebral que caracteriza esta etapa. Investigaciones recientes de \cite{swider-cios_young_2023} documentan que el uso de medios basados en pantalla durante estos años críticos puede influir de manera bidireccional en el desarrollo socioemocional: puede ser estimulante si existe un co-uso interactivo, pero resulta profundamente perjudicial si se utiliza como un "chupete electrónico" pasivo. En etapas posteriores, la plasticidad cerebral permite un aprendizaje rápido de interfaces complejas, pero al mismo tiempo facilita la formación de hábitos de consumo automatizados que erosionan las funciones ejecutivas y la capacidad de atención sostenida \cite{clemente-suarez_digital_2024}. \\

Un hito metodológico reciente lo constituye la revisión sistemática elaborada por \cite{page_psychological_2025}, la cual analiza la evolución de los daños psicológicos derivados de la Explotación y Abuso Sexual Infantil en Línea (OCSEA) a lo largo de más de tres décadas. Los autores concluyen que el daño emocional producto de la victimización digital es acumulativo, silencioso y a menudo invisible para el entorno físico del menor. Esta invisibilidad enfatiza que la prevención efectiva ya no puede depender exclusivamente de la observación conductual en el hogar, sino que requiere de herramientas computacionales capaces de identificar la victimización en sus etapas tempranas mediante el análisis directo de la conducta en la interfaz digital.\\

Desde una perspectiva de gobernanza macroestructural, el último reporte de la OECD \cite{noauthor_full_2025} revela una "vulnerabilidad estructural" inherente en las aplicaciones comerciales modernas. El informe denuncia la falta de estándares universales de seguridad por diseño (\textit{Safety by Design}), argumentando que el ecosistema actual prioriza la monetización de la atención infantil por encima de su integridad psicológica. Este documento fundamenta internacionalmente la necesidad de soluciones tecnológicas que no solo bloqueen contenidos mediante censura, sino que empoderen al menor y al núcleo familiar. La literatura académica actual converge en que la inteligencia artificial distribuida es la única herramienta técnica con la agilidad suficiente para seguir el ritmo de la constante evolución de las amenazas en línea, permitiendo una protección que es al mismo tiempo proactiva, local y contextualizada.\\

\section{Eficacia, usabilidad y vulnerabilidades de las herramientas actuales}

La evaluación técnica de las soluciones de control parental disponibles en el mercado revela brechas de seguridad sistémicas que ponen en duda su fiabilidad como mecanismos de protección primaria. En una rigurosa investigación publicada en \textit{IEEE Access}, \cite{rojas_technical_2025} analizaron las aplicaciones más populares de la Google Play Store y los servicios de filtrado de DNS Protegido (PDNS). Los resultados empíricos fueron alarmantes: una gran mayoría de estas herramientas fallaron catastróficamente en detectar sitios de alto riesgo cuando se ejecutaban en dispositivos no modificados (\textit{non-rooted}), y demostraron una vulnerabilidad crítica ante técnicas de evasión simples, como el uso de Redes Privadas Virtuales (VPN) gratuitas o navegadores web con DNS-over-HTTPS integrado. A este panorama se suma el desarrollo de arquitecturas como \textit{ChildProtect} \cite{ameur_childprotect_2023}, que intentan rastrear contenido hostil, pero que enfrentan el desafío endémico de mantener actualizadas las inmensas bases de datos de amenazas en tiempo real.\\



Más allá de la ineficacia técnica de bloqueo, la privacidad de los datos infantiles constituye un punto de preocupación académica severa. El estudio de \cite{feal_angel_2020} aplicó ingeniería inversa y descubrió que un porcentaje inaceptable de aplicaciones comerciales de control parental abusa rutinariamente de los permisos de accesibilidad de Android para recolectar información de uso, ubicación e interacciones, la cual es enviada a servidores de terceros y \textit{brokers} de datos con fines de segmentación publicitaria. Este hallazgo es reforzado por evaluaciones recientes de \cite{alqahtani_childrens_2023}, quienes advierten que muchas aplicaciones priorizan el rastreo geolocalizado (\textit{tracking}) sobre el filtrado inteligente de contenido. Por ello, la comunidad científica justifica que las nuevas propuestas deben procesar la información de manera local (\textit{Edge AI} / \textit{on-device}), garantizando criptográficamente que la bitácora de evidencias sea privada y esté protegida contra la extracción comercial.\\

Otro vector de fracaso ampliamente documentado en la literatura es la usabilidad (HCI) de estas herramientas. Investigaciones sobre la interacción humano-computadora en seguridad, como las de \cite{gnanasekaran_usability_2023}, señalan que los padres a menudo abandonan las aplicaciones de supervisión debido a interfaces de configuración excesivamente complejas y a la "fatiga de alerta" generada por una abrumadora cantidad de notificaciones irrelevantes. Los servicios basados exclusivamente en filtrado de red (DNS) suelen generar una alta tasa de falsos positivos, bloqueando recursos legítimos para tareas escolares mientras omiten contenido adulto reciente que aún no ha sido indexado. Para comprender mejor estas limitaciones, la siguiente tabla sintetiza la comparativa entre los enfoques actuales y las demandas de la investigación moderna:

\begin{table}[h!]
\centering
\caption{Comparativa de enfoques en herramientas de supervisión digital}
\begin{tabular}{|p{3cm}|p{4.5cm}|p{4.5cm}|p{2cm}|}
\hline
\textbf{Dimensión} & \textbf{Herramientas Tradicionales (Comerciales)} & \textbf{Enfoque Propuesto por la Literatura Actual} & \textbf{Referencia} \\ \hline
Método de Análisis & Filtrado de red (DNS, URLs) e historial de apps. & Análisis visual (\textit{Screen Capture}) semántico en tiempo real. & \cite{rojas_technical_2025} \\ \hline
Privacidad de Datos & Envío de datos a la nube, uso de rastreadores. & Procesamiento local (Edge AI), sin recolección en nube. & \cite{feal_angel_2020} \\ \hline
Usabilidad Parental & Paneles complejos, alta tasa de alertas falsas. & Bitácora de "momentos clave", alertas contextualizadas. & \cite{gnanasekaran_usability_2023} \\ \hline
Agencia del Menor & Bloqueo opaco, percepción de espionaje. & Transparencia, notificaciones de monitoreo educativo. & \cite{dumaru_its_2024} \\ \hline
\end{tabular}
\end{table}

\section{Detección automática de riesgos y el papel de la Inteligencia Artificial}

Ante las deficiencias de los filtros estáticos, la aplicación de técnicas de \textit{Machine Learning} y \textit{Deep Learning} ha demostrado resultados altamente prometedores en la detección de riesgos dentro de comunicaciones privadas cerradas. Investigaciones pioneras lideradas por \cite{razi_sliding_2023} analizaron corpus de mensajes directos (DMs) en Instagram, utilizando clasificadores avanzados basados en \textit{Random Forest} y modelos de Procesamiento de Lenguaje Natural (\gls{nlp}). Sus algoritmos lograron identificar episodios de acoso y solicitudes sexuales explícitas con una métrica de precisión (Área Bajo la Curva - $AUC$) de 0.88. Este enfoque revolucionario demuestra que la detección de riesgos no requiere una intrusión manual ni la lectura humana de los mensajes, sino que puede automatizarse mediante la identificación de patrones lingüísticos y secuencias conversacionales.\\



La transición de modelos simples de palabras clave hacia el análisis de intenciones es fundamental. Investigaciones centradas en la Interacción Humano-Computadora, como la de \cite{razi_human-centered_2021}, proponen el uso de redes neuronales recurrentes que analizan la secuencia histórica de los mensajes para detectar el \textit{grooming} en sus fases iniciales de "construcción de confianza". Los autores subrayan que los sistemas deben ser capaces de diferenciar el lenguaje coloquial, a menudo cargado de jerga y expresiones fuertes entre adolescentes, de las tácticas genuinas de ingeniería social utilizadas por depredadores adultos. La combinación de \gls{nlp} con la visión artificial local permite correlacionar el texto amenazante con las imágenes recibidas, generando una evaluación de riesgo multimodal mucho más precisa \cite{ali_getting_2023}.\\

Desde el punto de vista ético y algorítmico, \cite{agbana_ethical_2025} introducen la necesidad ineludible de la transparencia y la explicabilidad en los sistemas de \gls{ia} orientados a menores. Su investigación sugiere que el sistema nunca debe actuar como una "caja negra"; por el contrario, debe proporcionar evidencia clara al tutor de por qué intervino, manteniendo al humano en el ciclo de decisión (\textit{human-in-the-loop}). La implementación de una bitácora de evidencias criptográficamente protegida, como se plantea en esta tesis, sigue estrictamente estas recomendaciones, permitiendo que la \gls{ia} actúe como un filtro de triaje ultrarrápido y el padre actúe como el juez final del contexto. \\

La viabilidad práctica de ejecutar estos modelos de detección en tiempo real sobre teléfonos inteligentes se apoya fuertemente en investigaciones sobre arquitecturas de software ligeras. Según \cite{ali_efficient_2021} y \cite{sahu_edge_2025}, la aplicación de técnicas de optimización de borde (*Edge AI*) permite que las redes neuronales convolucionales analicen capturas de pantalla en milisegundos utilizando los aceleradores de hardware integrados del dispositivo (\gls{npu}/GPU). Esto no solo captura el riesgo en el instante exacto en que aparece en la pantalla del usuario (incluso en videos en vivo o transmisiones de juegos), sino que lo hace sin drenar la batería del dispositivo ni comprometer la fluidez de la interfaz gráfica, resolviendo así la principal barrera técnica para la adopción de este tipo de sistemas \cite{cao_mobile_2023}.\\

\section{Dinámicas de mediación parental y co-regulación tecnológica}

La eficacia de cualquier herramienta tecnológica de supervisión está inexorablemente vinculada al entorno sociológico del hogar y al estilo de crianza. La literatura sobre la sociología de la tecnología identifica dos estrategias principales: la \textit{mediación restrictiva} y la \textit{mediación habilitadora} \cite{livingstone_maximizing_2017}. Mientras la primera se basa en reglas estrictas, bloqueos y prohibiciones (reduciendo drásticamente la exposición a riesgos pero también limitando las oportunidades de aprendizaje y el desarrollo de habilidades digitales), la segunda fomenta el diálogo, el co-uso de los dispositivos y la construcción de criterio propio en el menor. Una revisión sistemática reciente realizada por \cite{banic_comparison_2024} destaca que la eficacia técnica de una aplicación de control parental se anula si no está respaldada por una relación de confianza; las herramientas deben servir como soporte para la comunicación y no como sustitutos de la crianza activa.\\



El diseño de las plataformas de seguridad debe, imperativamente, considerar la agencia y la voz del niño. En investigaciones seminales de co-diseño publicadas en foros como ACM CHI, \cite{mcnally_co-designing_2018} demostraron que los menores están dispuestos a aceptar la supervisión técnica siempre que esta sea transparente, predecible y posea un enfoque pedagógico. Este paradigma es respaldado por estudios recientes de \cite{dumaru_its_2024}, quienes evidenciaron que cuando los adolescentes entienden qué datos se están monitorizando y pueden acceder a sus propios reportes de actividad, desarrollan una mayor autorregulación y perciben a la aplicación como una herramienta de empoderamiento personal en lugar de un software de espionaje punitivo.\\

La necesidad de esta mediación asistida se vuelve crítica al analizar los factores de adicción. Investigaciones cualitativas mediante estudios de diario de \cite{theopilus_parental_2025} refuerzan que la mediación parental pasiva o negligente influye directamente en el riesgo de desarrollar Uso Problemático de Internet (UPI). En contraste, una mediación equilibrada y constante reduce la probabilidad de conductas compulsivas asociadas al consumo de redes sociales. Asimismo, \cite{allison_parental_2025} enfatizan que en la adolescencia temprana, los padres suelen recurrir a la restricción absoluta movidos por el pánico moral, pero advierten que la supervisión colaborativa es la única vía sostenible para desarrollar resiliencia digital a largo plazo. \\

La integración de la bitácora de evidencias asistida por \gls{ia} propuesta en esta tesis responde a este vacío teórico. Al proporcionar un registro visual y textual concreto de los incidentes, la aplicación otorga a los padres el "insumo" necesario para iniciar una conversación difícil (por ejemplo, sobre ciberacoso o contacto inapropiado) basándose en hechos objetivos y no en sospechas infundadas. Esta transparencia técnica facilita la transición de un modelo de "policía digital" a un modelo de "mentoría digital", transformando incidentes de riesgo aislados en lecciones valiosas que refuerzan el escudo cognitivo del menor frente a futuras amenazas.\\

\section{Epidemiología y prevalencia de riesgos críticos en el entorno móvil}

La justificación para implementar sistemas avanzados de análisis de pantalla radica en la abrumadora prevalencia epidemiológica de los riesgos digitales documentados por la comunidad científica. Un meta-análisis global liderado por \cite{mcarthur_global_2022} evidencia que el tiempo de uso recreativo de pantallas se ha disparado a niveles sin precedentes en la población escolar, estableciendo el sedentarismo digital como un vector principal de problemas de salud pública. Este incremento masivo de la exposición multiplica la superficie de ataque para los depredadores y la probabilidad estadística de encuentros hostiles en línea. \\

El ciberacoso sigue siendo la amenaza interpersonal más documentada. La revisión sistemática de \cite{zhu_cyberbullying_2021} reportó tasas de victimización por ciberacoso que alcanzan hasta el 57.5\% en diversas cohortes globales. El estudio identifica que el anonimato estructural de las plataformas y la desinhibición digital son los principales catalizadores de esta violencia verbal sostenida. En el ámbito escolar, investigaciones como las de \cite{kopecky_behaviour_2021} documentan que el uso ubicuo y no regulado de teléfonos móviles en los centros educativos facilita que las agresiones ocurridas en el aula continúen de manera incesante en el ciberespacio durante la noche, eliminando el "santuario" del hogar para las víctimas. \\



En cuanto a los riesgos relacionados con la sexualidad y el contenido explícito, estudios recientes en regiones en desarrollo, como el de \cite{yusuf_lets_2023}, arrojaron que el 17.4\% de los niños fueron expuestos de manera involuntaria a pornografía dura mientras navegaban aparentemente en sitios seguros. A la par, el fenómeno del \textit{sexting} y la extorsión asociada a la compartición de imágenes íntimas sin consentimiento (\textit{revenge porn}) se ha normalizado peligrosamente entre los adolescentes más jóvenes \cite{tomczyk_sexting_2023}. Esta normalización exige que las herramientas de supervisión posean capacidades de detección visual profunda capaces de identificar contenido desnudo o sugerente antes de que sea enviado o almacenado de forma permanente.\\

Finalmente, el impacto del exceso algorítmico y el tiempo de uso intensivo muestra una correlación irrefutable con la morbilidad psiquiátrica. \cite{bohnert_emerging_2021} analizaron datos de 13,203 niños, concluyendo que superar sistemáticamente las 3 horas diarias de uso interactivo correlaciona con declives severos en el bienestar socioemocional, generando cuadros de ansiedad y depresión clínica \cite{marciano_digital_2022}. Esta vulnerabilidad psicológica actúa como un bucle de retroalimentación fatal: según \cite{el-asam_psychological_2022}, el malestar psicológico previo (como la soledad o la marginación) actúa como el principal mediador del riesgo en línea, empujando a los niños vulnerables a buscar validación en extraños y exponiéndose al \textit{grooming}. Esto valida que las herramientas de protección móvil deben funcionar no solo como escudos técnicos, sino como termómetros de salud mental preventiva.

\section{Impacto de las arquitecturas persuasivas y algoritmos de recomendación}

Una línea de investigación crítica que ha ganado tracción recientemente es el estudio de cómo la arquitectura misma de las plataformas digitales induce comportamientos de riesgo sin la intervención de un agente humano malicioso (como un acosador). Los menores interactúan diariamente con sistemas diseñados mediante "algoritmos de recomendación" cuya única métrica de éxito es la retención de la atención (\textit{engagement}). Según el análisis de \cite{selak_problematic_2025}, este diseño persuasivo genera dinámicas de Uso Problemático de Internet, donde el menor experimenta una pérdida total de la noción del tiempo y de su capacidad de elección consciente, cayendo en bucles de visualización pasiva (\textit{doomscrolling}).\\

La gravedad de estos algoritmos reside en su capacidad para crear "cámaras de eco" o "burbujas de filtro". Un menor que visualiza por curiosidad un video relacionado con trastornos alimenticios o autolesiones puede ver su \textit{feed} rápidamente inundado por contenido similar, empujándolo hacia nichos de radicalización o daño autoinfligido. La comunidad académica, como señala \cite{chou_rectifying_2023}, está investigando activamente métodos de "rectificación algorítmica" y mitigación de sesgos, argumentando que las plataformas tienen una responsabilidad fiduciaria de no servir contenido extremo a mentes en desarrollo.\\

Dado que las corporaciones tecnológicas son reacias a modificar sus lucrativos algoritmos, la protección técnica recae sobre la supervisión en el dispositivo del usuario final. Los filtros de red clásicos son ciegos ante el flujo de contenido generado dinámicamente dentro de la aplicación de TikTok, Instagram o YouTube Shorts. Solo un sistema basado en \gls{ia} que capture y procese la pantalla en milisegundos puede evaluar si el algoritmo de la plataforma está sometiendo al menor a un bombardeo algorítmico tóxico, justificando así la necesidad imperativa de evolucionar hacia la inspección visual local propuesta a lo largo de este estado del arte.\\